% ============================================================
\chapter{Mesures de Risque}
\label{ch:risque}
% ============================================================

\section{Introduction}

Apr\`es les crises financi\`eres des ann\'ees 1990 --- d\'eb\^acle d'Orange County (1994), faillite de Barings (1995), crise LTCM (1998) --- la question de la mesure du risque est devenue centrale. J.P.\,Morgan popularise la \emph{Value at Risk} (VaR) en 1994 avec son syst\`eme RiskMetrics~: un chiffre unique r\'esumant la perte maximale probable sur un horizon donn\'e avec un certain niveau de confiance. Mais Philippe Artzner et ses coll\`egues montrent en 1999 que la VaR n'est pas \emph{coh\'erente}~: elle ne respecte pas la sous-additivit\'e, ce qui signifie que diversifier un portefeuille peut, paradoxalement, augmenter la VaR mesur\'ee. L'\emph{Expected Shortfall} (CVaR) corrige ce d\'efaut. Ce chapitre parcourt cette histoire, des mesures empiriques aux axiomes th\'eoriques, en passant par les copules et les tests de r\'esistance.

\begin{intuition}
Une mesure de risque r\'epond \`a la question : \og Combien puis-je perdre au
maximum avec une probabilit\'e donn\'ee sur un horizon donn\'e ? \fg La VaR
r\'epond par un quantile, mais ne dit rien sur l'ampleur des pertes au-del\`a.
L'Expected Shortfall comble cette lacune en moyennant les pertes dans la queue
de la distribution.
\end{intuition}

\section{Value at Risk (VaR)}

\begin{definition}[Value at Risk]
\index{Value at Risk}
Soit $L$ la variable al\'eatoire repr\'esentant la perte d'un portefeuille sur
un horizon $\Delta t$. La \textbf{Value at Risk}\index{VaR} au niveau de
confiance $\alpha \in (0, 1)$ est
\[
    \mathrm{VaR}_\alpha(L) = \inf\{l \in \R : \PP(L > l) \leq 1 - \alpha\}
    = F_L^{-1}(\alpha),
\]
o\`u $F_L^{-1}$ d\'esigne le quantile g\'en\'eralis\'e de $L$.
\end{definition}

\begin{exemple}[VaR gaussienne]
Si $L \sim \mathcal{N}(\mu, \sigma^2)$, alors
$\mathrm{VaR}_\alpha = \mu + \sigma\,\Phi^{-1}(\alpha)$,
o\`u $\Phi$ est la fonction de r\'epartition de la loi normale standard.
Pour $\alpha = 0{,}99$ : $\mathrm{VaR}_{99\%} = \mu + 2{,}326\,\sigma$.
\end{exemple}

\begin{proposition}[M\'ethodes de calcul de la VaR]
\begin{enumerate}
    \item \textbf{Historique} : $\mathrm{VaR}_\alpha = $ quantile empirique
    des rendements pass\'es.
    \item \textbf{Param\'etrique (variance-covariance)} : suppose une
    distribution (souvent gaussienne) et utilise $\mu$, $\sigma$ estim\'es.
    \item \textbf{Monte Carlo} : simulation de sc\'enarios et calcul du
    quantile sur les pertes simul\'ees.
\end{enumerate}
\end{proposition}

\begin{attention}[title=\textbf{Limites de la VaR}]
La VaR ne satisfait pas la propri\'et\'e de sous-additivit\'e en g\'en\'eral :
on peut avoir $\mathrm{VaR}_\alpha(L_1 + L_2) > \mathrm{VaR}_\alpha(L_1) +
\mathrm{VaR}_\alpha(L_2)$. Cela signifie que la diversification peut
\emph{sembler} augmenter le risque, ce qui est un d\'efaut conceptuel majeur.
\end{attention}

\section{Expected Shortfall (CVaR)}

\begin{definition}[Expected Shortfall]
\index{Expected Shortfall}\index{CVaR}
L'\textbf{Expected Shortfall} (ou CVaR, Conditional Value at Risk) au niveau
$\alpha$ est
\[
    \mathrm{ES}_\alpha(L) = \frac{1}{1-\alpha}\int_\alpha^1
    \mathrm{VaR}_u(L)\,\mathrm{d}u.
\]
Si $F_L$ est continue, $\mathrm{ES}_\alpha(L) = \E[L \mid L \geq \mathrm{VaR}_\alpha(L)]$.
\end{definition}

\begin{exemple}[ES gaussien]
Si $L \sim \mathcal{N}(\mu, \sigma^2)$ :
\[
    \mathrm{ES}_\alpha = \mu + \sigma\,\frac{\phi(\Phi^{-1}(\alpha))}{1 - \alpha},
\]
o\`u $\phi$ est la densit\'e de la loi normale standard. Pour $\alpha = 0{,}99$ :
$\mathrm{ES}_{99\%} \approx \mu + 2{,}665\,\sigma$.
\end{exemple}

\begin{theoreme}[Repr\'esentation duale de l'ES]
\index{repr\'esentation duale}
\[
    \mathrm{ES}_\alpha(L) = \sup_{Q \in \mathcal{Q}_\alpha} \E_Q[L],
\]
o\`u $\mathcal{Q}_\alpha = \{Q \ll \PP : \mathrm{d}Q/\mathrm{d}\PP \leq 1/(1-\alpha)\}$.
Cette repr\'esentation montre que l'ES est le pire cas d'esp\'erance sur un
ensemble d'ambig\"uit\'e.
\end{theoreme}

\section{Mesures de risque coh\'erentes}

\begin{definition}[Mesure de risque coh\'erente]
\index{mesure de risque coh\'erente}
Une fonctionnelle $\rho : \mathcal{X} \to \R$ est une \textbf{mesure de risque
coh\'erente} (Artzner, Delbaen, Eber, Heath, 1999) si elle satisfait :
\begin{enumerate}
    \item \textbf{Monotonie} : $L_1 \leq L_2$ p.s. $\implies$ $\rho(L_1) \leq \rho(L_2)$.
    \item \textbf{Invariance par translation} : $\rho(L + c) = \rho(L) + c$
    pour tout $c \in \R$.
    \item \textbf{Homog\'en\'eit\'e positive} : $\rho(\lambda L) = \lambda\rho(L)$
    pour tout $\lambda > 0$.
    \item \textbf{Sous-additivit\'e} : $\rho(L_1 + L_2) \leq \rho(L_1) + \rho(L_2)$.
\end{enumerate}
\end{definition}

\begin{theoreme}[ADEH, 1999]
\begin{enumerate}
    \item L'Expected Shortfall est une mesure de risque \textbf{coh\'erente}.
    \item La VaR n'est \textbf{pas} coh\'erente (elle viole la sous-additivit\'e).
    \item Toute mesure de risque coh\'erente admet une repr\'esentation
    $\rho(L) = \sup_{Q \in \mathcal{Q}} \E_Q[L]$ pour un ensemble convexe
    ferm\'e $\mathcal{Q}$ de mesures de probabilit\'e.
\end{enumerate}
\end{theoreme}

\begin{formules}[title=\textbf{Comparaison VaR vs ES}]
\begin{center}
\begin{tabular}{lcc}
\textbf{Propri\'et\'e} & \textbf{VaR} & \textbf{ES} \\
\hline
Sous-additivit\'e & Non & Oui \\
Coh\'erence & Non & Oui \\
Sensibilit\'e \`a la queue & Non & Oui \\
R\'eglementation B\^ale III+ & Non & Oui \\
Facilit\'e de backtesting & Oui & Difficile \\
\end{tabular}
\end{center}
\end{formules}

\section{Tests de r\'esistance (Stress Testing)}

\begin{definition}[Stress test]
\index{stress test}
Un \textbf{test de r\'esistance} consiste \`a \'evaluer les pertes d'un
portefeuille sous des sc\'enarios extr\^emes pr\'ed\'efinis ou historiques :
\[
    L_{\text{stress}} = \sum_{i=1}^n \Delta_i \cdot S_i^{\text{stress}},
\]
o\`u $\Delta_i$ sont les sensibilit\'es et $S_i^{\text{stress}}$ les chocs
appliqu\'es aux facteurs de risque.
\end{definition}

\begin{proposition}[Types de stress tests]
\begin{enumerate}
    \item \textbf{Historiques} : reproduction de crises pass\'ees (2008, 2020).
    \item \textbf{Hypoth\'etiques} : sc\'enarios construits (choc de taux de
    300 bp, effondrement d'un secteur).
    \item \textbf{Reverse stress tests} : recherche du sc\'enario qui m\`ene
    \`a une perte critique donn\'ee.
\end{enumerate}
\end{proposition}

\section{Copules et mod\'elisation de la d\'ependance}

\begin{definition}[Copule]
\index{copule}
Une \textbf{copule} $C : [0,1]^d \to [0,1]$ est une fonction de r\'epartition
dont les marginales sont uniformes sur $[0,1]$. Elle capture la structure de
d\'ependance ind\'ependamment des lois marginales.
\end{definition}

\begin{theoreme}[Sklar, 1959]
\index{th\'eor\`eme de Sklar}
Pour toute loi jointe $F$ de marginales $F_1, \ldots, F_d$, il existe une
copule $C$ telle que
\[
    F(x_1, \ldots, x_d) = C(F_1(x_1), \ldots, F_d(x_d)).
\]
Si les $F_i$ sont continues, $C$ est unique.
\end{theoreme}

\begin{exemple}[Copule gaussienne]
\index{copule!gaussienne}
La \textbf{copule gaussienne} de param\`etre $\Sigma$ (matrice de corr\'elation) est
\[
    C_\Sigma^{\mathrm{Ga}}(u_1, \ldots, u_d)
    = \Phi_\Sigma(\Phi^{-1}(u_1), \ldots, \Phi^{-1}(u_d)),
\]
o\`u $\Phi_\Sigma$ est la fonction de r\'epartition de la loi normale
multidimensionnelle.
\end{exemple}

\begin{definition}[Copule de Clayton]
\index{copule!Clayton}
La \textbf{copule de Clayton} (bivari\'ee) est
\[
    C_\theta(u, v) = \left(u^{-\theta} + v^{-\theta} - 1\right)^{-1/\theta},
    \quad \theta > 0.
\]
Elle mod\'elise une forte d\'ependance dans la queue inf\'erieure (risque de
co-crash).
\end{definition}

\begin{remarque}
La copule gaussienne sous-estime la d\'ependance de queue. C'est l'une des
causes de la crise des subprimes de 2008, o\`u les mod\`eles de CDO
utilisaient la copule gaussienne de Li (2000) sans tenir compte de la
d\'ependance extr\^eme.
\end{remarque}

\begin{definition}[D\'ependance de queue]
\index{d\'ependance de queue}
Le \textbf{coefficient de d\'ependance de queue inf\'erieure} est
\[
    \lambda_L = \lim_{u \to 0^+} \frac{C(u, u)}{u}.
\]
Pour la copule gaussienne, $\lambda_L = 0$ (sauf si $\rho = 1$).
Pour la copule de Clayton, $\lambda_L = 2^{-1/\theta} > 0$.
\end{definition}

\begin{codeblock}[title=\textbf{Calcul de la VaR et de l'ES par simulation Monte Carlo}]
\begin{minted}{python}
import numpy as np
from scipy.stats import norm

np.random.seed(42)
n_simulations = 100_000
mu, sigma = 0.0, 0.02  # rendement journalier
alpha = 0.99

# Simulation des rendements
returns = np.random.normal(mu, sigma, n_simulations)
losses = -returns  # perte = -rendement

# VaR historique
var_99 = np.percentile(losses, 100 * alpha)

# Expected Shortfall
es_99 = losses[losses >= var_99].mean()

print(f"VaR 99%: {var_99:.4f}")
print(f"ES 99%:  {es_99:.4f}")
print(f"VaR analytique: {-mu + sigma * norm.ppf(alpha):.4f}")
\end{minted}
\end{codeblock}

\section{Exercices}

\begin{exercice}[VaR d'un portefeuille]
Un portefeuille de valeur 10 M\texteuro{} a des rendements journaliers
$\mathcal{N}(0{,}05\%,\, (1{,}5\%)^2)$. Calculer la VaR \`a 99\% \`a 1 jour
et \`a 10 jours (r\`egle de la racine carr\'ee du temps).
\end{exercice}

\begin{exercice}[Violation de la sous-additivit\'e]
Construire un exemple explicite de deux variables al\'eatoires $L_1, L_2$
telles que $\mathrm{VaR}_{0{,}95}(L_1 + L_2) > \mathrm{VaR}_{0{,}95}(L_1) +
\mathrm{VaR}_{0{,}95}(L_2)$.
\textit{Indication} : utiliser des pertes de Bernoulli concentr\'ees.
\end{exercice}

\begin{exercice}[Copule de Clayton]
Simuler 10\,000 paires $(U, V)$ suivant une copule de Clayton avec
$\theta = 2$. Estimer le coefficient de d\'ependance de queue inf\'erieure
et comparer avec la valeur th\'eorique $\lambda_L = 2^{-1/2}$.
\end{exercice}

\begin{exercice}[Mesure de risque convexe]
Montrer que toute mesure de risque coh\'erente est convexe. Donner un exemple
de mesure convexe qui n'est pas coh\'erente (homog\'en\'eit\'e positive
viol\'ee).
\end{exercice}

\begin{exercice}[Stress test]
Concevoir un stress test historique bas\'e sur la crise de 2008 pour un
portefeuille contenant 60\% d'actions et 40\% d'obligations. Calculer la
perte sous ce sc\'enario et comparer avec la VaR \`a 99\%.
\end{exercice}

\begin{exercice}[Backtesting]
Expliquer le test de Kupiec pour le backtesting de la VaR. Si sur 250 jours
de trading, on observe 5 d\'epassements de la VaR \`a 99\%, la VaR est-elle
correctement calibr\'ee au seuil de 5\% ?
\end{exercice}
